\uuid{NJBP}
\titre{DL en l'infini/Asymptotes}
\auteur{hanani ; mzari}
\organisation{exo7}
\licence{CC BY-NC-SA 4.0}
\datecreate{2020-01-20}
\niveau{L1}
\module{Analyse}
\chapitre{Développement limité}
\sousChapitre{Applications}
\difficulte{3}
\theme{Applications des DL, DL en l'infini, Asymptotes}
\exo7qcmid{lille-2.developpements-limites.038}

\contenu{
\texte{
Soit $a\in \Rr$ et $f$ une fonction. On dit que $f(x)$ est un $o\left((x-a)^n\right)$ au voisinage de $a$ et on écrit $f(x)=o((x-a)^n)$ si $\displaystyle \lim _{x\to a}\frac{f(x)}{(x-a)^n}=0$. On note par $DL_n(a)f(x)$ le développement limité de $f$ à l'ordre $n$ en $a$. On note aussi $DL_n(a^+)f(x)$ (resp. $DL_n(a^-)f(x)$) le développement limité de $f$ à l'ordre $n$ à droite en $a$ (resp. à gauche en $a$).
}

\begin{question}
\qclassification{125.03}
\qtags{motcle=DL en l'infini/Asymptotes}
Soit $\displaystyle f(x)=\sqrt{2+x^2}$. On note $\Gamma$ le graphe de $f$ et on pose $\displaystyle g(x)=xf\left(\frac{1}{x}\right)$. Parmi les affirmations suivantes, cocher celles qui sont vraies :
\begin{answers}
\bad{$\displaystyle DL_2(0)g(x)=1+x+x^2+o(x^2)$.}
\good{Au voisinage de $+\infty$, on a : $\displaystyle f(x)=x+\frac{1}{x}+o\left(\frac{1}{x}\right)$.}
\good{$\Gamma$ admet la droite d'équation $y=x$ comme asymptote au voisinage de $+\infty$.}
\bad{$\Gamma$ est en dessous de la droite d'équation $y=x$ au voisinage de $+\infty$.}
\end{answers}
\begin{explanations}
D'abord, au voisinage de $0^+$, on a : $g(x)=\sqrt{1+2x^2}=1+x^2+o(x^2)$ et donc, au voisinage de $+\infty$, on a :
$$f(x)=x+\frac{1}{x}+o\left(\frac{1}{x}\right)\mbox{ car }f(x)=xg\left(\frac{1}{x}\right).$$
La droite d'équation $y=x$ est une asymptote au voisinage de $+\infty$. De plus, au voisinage de $+\infty$, $\displaystyle f(x)-y=\frac{1}{x}+o\left(\frac{1}{x}\right)\geq 0$. Donc $\Gamma$ est au dessus de la droite d'équation $y=x$ au voisinage de $+\infty$.
\end{explanations}
\end{question}
}
